% =====================================================================
\section{Características estáticas}
\label{sec:caracteristicas}
% =====================================================================

\subsection{Curva de entrada (emisor común)}

La entrada en emisor común es la juntura E--B: la característica $I_B$ vs.\ $V_{EB}$ (con
$V_{EC}$ fija) es la de un diodo \diapo{1.28}. Como $I_B=I_C/\beta$, la curva $I_C$ vs.\
$V_{EB}$ (característica de \emph{transferencia}) tiene la misma forma, escalada por $\beta$.
En Si conduce apreciablemente a partir de \SIrange{0.5}{0.6}{V}, y a corriente constante
$V_{EB}$ \textbf{disminuye unos \SI{2}{mV/\celsius}} con la temperatura (la curva se corre a la
izquierda al calentarse, porque $n_i^2$ crece).

\begin{ojo}[\diapo{1.28}]
Las curvas mostradas están rotuladas $i_C$ vs.\ $v_{BE}$: son curvas de transferencia, no de
entrada. La forma es la misma (exponencial), así que la conclusión ``es la curva de un
diodo'' vale; pero la corriente de entrada es $i_B$, $\beta$ veces menor.
\end{ojo}

\subsection{Curvas de salida (emisor común)}

Las curvas de salida grafican $I_C$ en función de $V_{EC}$, una curva por cada valor de $I_B$
\diapo{1.29}. Las de la figura~\ref{fig:salida} están calculadas con Ebers--Moll y los
parámetros del ejemplo, agregando el efecto Early.

\begin{figure}[H]
\centering
\begin{tikzpicture}
\begin{axis}[
  width=0.82\textwidth,height=7cm,
  xmin=0,xmax=10,ymin=0,ymax=2.6,
  xlabel={$V_{EC}$ (V)},ylabel={$I_C$ (mA)},
  axis lines=left,grid=major,grid style={gray!20},
  every axis plot/.append style={line width=1.2pt,azul,no markers},
  clip=false,
]
\fill[naranjaclaro,opacity=0.7] (axis cs:0,0) rectangle (axis cs:0.25,2.6);
\fill[gris!15] (axis cs:0,0) rectangle (axis cs:10,0.04);
\addplot table[x=VEC,y=IC1] {datos/salida.dat};
\addplot table[x=VEC,y=IC2] {datos/salida.dat};
\addplot table[x=VEC,y=IC3] {datos/salida.dat};
\addplot table[x=VEC,y=IC4] {datos/salida.dat};
\node[font=\scriptsize,anchor=west] at (axis cs:10,0.60) {$I_B=\SI{1}{\micro A}$};
\node[font=\scriptsize,anchor=west] at (axis cs:10,1.20) {\SI{2}{\micro A}};
\node[font=\scriptsize,anchor=west] at (axis cs:10,1.80) {\SI{3}{\micro A}};
\node[font=\scriptsize,anchor=west] at (axis cs:10,2.40) {\SI{4}{\micro A}};
\node[naranja!70!black,font=\small,anchor=west] at (axis cs:0.3,0.25) {$\leftarrow$ saturación};
\node[font=\small] at (axis cs:5,2.05) {activa directa};
\node[gris,font=\scriptsize,anchor=south west] at (axis cs:3,0.04) {corte ($I_B=0$)};
\end{axis}
\end{tikzpicture}
\caption{Curvas de salida de un PNP en emisor común (calculadas). En activa las curvas están
equiespaciadas ($\Delta I_C=\beta\Delta I_B$) y levemente inclinadas (efecto Early). Por
debajo de $V_{EC}\approx\SIrange{0.1}{0.2}{V}$ la juntura C--B entra en directa y el
transistor se satura: $I_C$ cae rápidamente a cero.}
\label{fig:salida}
\end{figure}

\paragraph{Base común.} Si se grafica $I_C$ vs.\ $V_{CB}$ con $I_E$ como parámetro, las curvas
son prácticamente horizontales ($I_C\approx\adc I_E$), sin efecto Early visible, y siguen
así hasta \emph{pasar} $V_{CB}=0$: la saturación empieza recién cuando la juntura C--B entra
claramente en directa ($V_{CB}\approx\SI{0.5}{V}$).

\subsection{Corrientes de fuga: $I_{CB0}$ e $I_{CE0}$}

Incluyendo la corriente inversa de la juntura C--B \diapo{2.9--2.11}:

\begin{clave}[Base común y emisor común \enhoja]
\[
  I_C = \adc I_E + I_{CB0}
  \qquad\Longrightarrow\qquad
  I_C = \frac{\adc}{1-\adc}I_B + \frac{I_{CB0}}{1-\adc} = \bdc I_B + I_{CE0},
\]
\[
  I_{CE0} = \frac{I_{CB0}}{1-\adc} = (\bdc+1)\,I_{CB0}.
\]
\begin{itemize}[nosep]
  \item $I_{CB0}$: corriente de colector con el \textbf{emisor abierto} ($I_E=0$).
  \item $I_{CE0}$: corriente de colector con la \textbf{base abierta} ($I_B=0$).
\end{itemize}
\end{clave}

\begin{intuicion}[por qué $I_{CE0}\gg I_{CB0}$]
Con la base abierta, los electrones de la corriente de fuga C--B que llegan a la base no
tienen por dónde salir: se acumulan, polarizan en directa a la juntura E--B y el emisor
inyecta huecos. Esa fuga actúa como una corriente de base y el transistor la
\emph{amplifica} por $\beta+1$. Por eso, y porque $I_{CB0}$ se duplica cada \SI{10}{\celsius}
aproximadamente, la fuga es un problema en circuitos analógicos que hay que compensar
\diapo{2.34}.
\end{intuicion}

\begin{ojo}[$I_{CB0}$ \diapo{2.11}]
``En nuestro desarrollo sería la $I_{Cn}$'': es una aproximación. Con Ebers--Moll, emisor
abierto ($I_E=0$) y $V_{CB}\ll-V_T$:
\[
  I_{CB0} = I_{R0}\left(1-\aF\aR\right),
\]
que incluye $I_{Cn}$ más la pequeña corriente de huecos de la base hacia el colector. Además
hay que sumarle la corriente de generación en la ZCE de C--B, que en Si suele dominar.
\end{ojo}

\subsection{Gummel plot y dependencia de $\beta$}

$\beta$ es la figura de mérito más usada para comparar transistores, pero \textbf{no es
constante}: depende de $I_C$ y de la temperatura \diapo{1.27}. El \emph{gráfico de
Gummel} muestra $\log I_C$ y $\log I_B$ en función de $V_{EB}$; la distancia vertical entre
ambas curvas es $\log\beta$.

\begin{figure}[H]
\centering
\begin{tikzpicture}
\begin{semilogyaxis}[
  name=gum,
  width=0.47\textwidth,height=6.2cm,
  xmin=0.2,xmax=0.95,ymin=1e-13,ymax=1e-1,
  xlabel={$V_{EB}$ (V)},ylabel={corriente (A)},
  legend pos=north west,legend style={font=\scriptsize},
  every axis plot/.append style={line width=1.1pt},
  ytick={1e-12,1e-10,1e-8,1e-6,1e-4,1e-2},
]
\addplot[azul] table[x=VEB,y=IC] {datos/gummel.dat};
\addlegendentry{$I_C$}
\addplot[rojo] table[x=VEB,y=IB] {datos/gummel.dat};
\addlegendentry{$I_B$}
\draw[<->,gris] (axis cs:0.6,2.2e-4) -- node[right,font=\scriptsize,black] {$\beta$} (axis cs:0.6,7e-7);
\node[font=\scriptsize,align=center,rojo] at (axis cs:0.36,2e-8) {R--G en la\\ZCE E--B};
\node[font=\scriptsize,align=center,azul] at (axis cs:0.83,1e-6) {alta\\inyección};
\end{semilogyaxis}
\begin{axis}[
  at={(gum.east)},anchor=west,xshift=1.4cm,
  width=0.47\textwidth,height=6.2cm,
  xmode=log,xmin=1e-6,xmax=1e-1,ymin=0,ymax=360,
  xlabel={$I_C$ (A)},ylabel={$\beta$},
  every axis plot/.append style={line width=1.1pt,smooth},
  legend pos=north west,legend style={font=\scriptsize},
]
\addplot[rojo] coordinates {(1e-6,150)(1e-5,215)(1e-4,270)(1e-3,305)(1e-2,275)(3e-2,200)(1e-1,95)};
\addlegendentry{\SI{125}{\celsius}}
\addplot[azul] coordinates {(1e-6,100)(1e-5,145)(1e-4,180)(1e-3,195)(1e-2,170)(3e-2,120)(1e-1,55)};
\addlegendentry{\SI{25}{\celsius}}
\addplot[violeta] coordinates {(1e-6,65)(1e-5,90)(1e-4,105)(1e-3,110)(1e-2,90)(3e-2,60)(1e-1,25)};
\addlegendentry{\SI{-55}{\celsius}}
\end{axis}
\end{tikzpicture}
\caption{Izquierda: gráfico de Gummel (Pierret, fig.~11.13). Derecha: $\beta$ en función de la
corriente de colector y la temperatura (valores típicos, \diapo{1.27}). El modelo ideal sólo
vale en el tramo central donde las dos rectas son paralelas.}
\label{fig:gummel}
\end{figure}

\begin{itemize}
  \item \textbf{Corrientes bajas:} a $I_B$ se le suma la corriente de recombinación en la ZCE
  de E--B, que crece como $e^{V_{EB}/2V_T}$ (pendiente mitad). $I_B$ ya no es proporcional a
  $I_C$ y $\beta$ cae.
  \item \textbf{Corrientes altas:} alta inyección en la base (la pendiente de $I_C$ pasa a
  $e^{V_{EB}/2V_T}$), aglomeración de corriente en los bordes del emisor
  (\emph{current crowding}) y caídas en las resistencias serie. $\beta$ vuelve a caer.
  \item \textbf{Temperatura:} $\beta$ aumenta con $T$ (sobre todo porque mejora la eficiencia
  de emisión).
\end{itemize}

Por eso, al decir ``el transistor tiene $\beta=200$'' hay que aclarar a qué corriente y
temperatura. En la materia se toma $\beta$ constante dentro del rango de validez del modelo
\diapo{1.31}.
